\chapter{Wariacje struktur Hodge'a, koneksja Gaussa--Manina i degeneracje}
\label{ch:xv-vhs-degeneracje}

\section{Dlaczego struktura Hodge'a musi się zmieniać?}
W rozdziale~\ref{ch:xiv-hodge-monodromia} wyposażyliśmy
$H^k(X;\Q)$ gładkiej odmiany projektowej $X$ w strukturę Hodge'a
oraz obliczyliśmy monodromię izolowanej osobliwości kwadratowej.
Teraz rozważymy \emph{rodzinę} $f:\mathcal X\to S$ odmian i zapytamy,
jak przy zmianie parametru $s\in S$ porównywać podprzestrzenie
$H^{p,q}(X_s)\subset H^k(X_s;\C)$.
Grupy kohomologii są lokalnie identyczne jako przestrzenie wektorowe,
lecz filtracja Hodge'a może się zmieniać. To rozróżnienie prowadzi
do koneksji Gaussa--Manina, równań Picarda--Fuchsa i granicznych
struktur mieszanych.

\begin{definicja}[Gładka właściwa rodzina projektowa]
Niech $S$ będzie spójną gładką odmianą algebraiczną nad $\C$
(rozpatrywaną także jako rozmaitość analityczna),
$\mathcal X$ gładką odmianą algebraiczną nad $\C$,
a $f:\mathcal X\to S$ gładkim, surjektywnym morfizmem
projektowym. Odpowiadające mu odwzorowanie analityczne
jest właściwą submersją holomorficzną. Włóknem w $s\in S$ nazywamy
$X_s=f^{-1}(s)$. Każde $X_s$ jest wówczas zwartą gładką
odmianą projektową. W tym rozdziale właśnie taką rodzinę
nazywamy krótko \emph{gładką rodziną projektową}.
\end{definicja}
\begin{uwaga}[Po co wszystkie hipotezy?]
Właściwość i warunek submersji pozwalają zastosować twierdzenie
Ehresmanna, a projektowość zapewnia rozkład Hodge'a oraz
jego polaryzację. Na bazie $S$ nie musi istnieć globalny wybór
baz kohomologii włókien: przeszkodą jest monodromia.
\end{uwaga}

\section{Lokalny system, płaska wiązka i koneksja Gaussa--Manina}
\begin{definicja}[Lokalny system kohomologii]
Niech $f:\mathcal X\to S$ będzie gładką rodziną projektową,
$k\ge0$. Oznaczmy przez
$\mathbb V_{\Z}:=R^kf_*\Z_{\mathcal X}/\mathrm{tors}$
lokalny system wolnych grup abelowych skończonego rzędu,
przez $\mathbb V_{\C}=\mathbb V_{\Z}\otimes_{\Z}\C$
jego zespolenie, a przez
$\mathcal H:=\mathbb V_{\C}\otimes_{\C_S}\mathcal O_S$
skończonego rzędu holomorficzną wiązkę wektorową
(tutaj $\C_S$ oznacza snop stały na $S$).
Przy przejściu od współczynników całkowitych do zespolonych
torsja znika, stąd naturalnie
$\mathbb V_{\C}\cong R^kf_*\C_{\mathcal X}$.
Dokładniej, iloraz przez torsję bierze się włókno po włóknie;
z twierdzenia Ehresmanna otrzymujemy lokalny system
$\mathbb V_{\Z}$ o włóknie $H^k(X_s;\Z)/\mathrm{tors}$.
\end{definicja}
\begin{stwierdzenie}[Płaska koneksja z lokalnego systemu]
\label{prop:xv-gm-local}
Niech $S$ będzie zespoloną rozmaitością, a
$\mathbb V$ lokalnym systemem skończeniewymiarowych
zespolonych przestrzeni wektorowych na $S$. Wówczas wiązka
$\mathcal H=\mathbb V\otimes_{\C_S}\mathcal O_S$
ma kanoniczną koneksję holomorficzną
$\nabla:\mathcal H\to\Omega_S^1\otimes_{\mathcal O_S}\mathcal H$
charakteryzowaną przez $\nabla(v\otimes g)=v\otimes\dd g$
dla lokalnie stałego $v$. Jest ona płaska i jej lokalnymi
przekrojami poziomymi są dokładnie przekroje $\mathbb V$.
\end{stwierdzenie}
\begin{proof}
Wybierzmy pokrycie $S$ zbiorami po prostu spójnymi $U_\alpha$,
na których $\mathbb V$ jest stały, i lokalne bazy
$e_{\alpha,1},\dots,e_{\alpha,r}$. Dla
$\sigma=\sum_ig_ie_{\alpha,i}$ przyjmujemy
$\nabla\sigma=\sum_i\dd g_i\otimes e_{\alpha,i}$.
Na przecięciu $U_\alpha\cap U_\beta$ macierz przejścia
między bazami ma lokalnie stałe współczynniki, a więc
jej różniczka znika. Zatem wzór nie zależy od wyboru bazy.
Reguła Leibniza wynika z $\dd(hg)=h\dd g+g\dd h$,
natomiast $\nabla^2=0$ z $\dd^2=0$ w płaskich bazach.
Warunek $\nabla\sigma=0$ oznacza $\dd g_i=0$ dla każdego $i$,
czyli lokalną stałość wszystkich $g_i$.
\end{proof}
\begin{definicja}[Koneksja Gaussa--Manina]
Niech $f:\mathcal X\to S$ będzie gładką
rodziną projektową, a $k\ge0$. Koneksję z
Stwierdzenia~\ref{prop:xv-gm-local}, zastosowaną do
lokalnego systemu $\mathbb V_{\C}=R^kf_*\C_{\mathcal X}$,
nazywamy \emph{koneksją Gaussa--Manina}
rodziny $f$ w stopniu $k$.
\end{definicja}
\begin{stwierdzenie}[Różniczkowanie okresów]
\label{prop:xv-period-derivative}
Niech $S$ będzie gładką zespoloną krzywą algebraiczną,
$f:\mathcal X\to S$ gładką rodziną projektową, $k\ge0$,
$\mathcal H=(R^kf_*\C_{\mathcal X})\otimes_{\C_S}\mathcal O_S$
wiązką z koneksją Gaussa--Manina,
a $U\subset S$ dyskiem współrzędnych z parametrem $t$.
Dla $t\in U$ oznaczmy $X_t=f^{-1}(t)$.
Niech $\gamma_t\in H_k(X_t;\Z)$ będzie
lokalnie poziomą rodziną cykli w sensie lokalnego
systemu homologii, a
$\omega\in\Gamma(U,\mathcal H)$ holomorficznym przekrojem.
Wówczas dla funkcji okresowej
$\Pi(t)=\langle\gamma_t,\omega(t)\rangle$
zachodzi
$\Pi'(t)=\langle\gamma_t,\nabla_{\partial/\partial t}\omega(t)\rangle$.
\end{stwierdzenie}
\begin{proof}
Na $U$ możemy jednocześnie wybrać płaskie bazy homologii
oraz kohomologii. W tych bazach $\gamma_t$ ma stałe
współrzędne, a $\omega(t)$ ma współrzędne zmienne.
Różniczkujemy zwykłą skończoną sumę iloczynów współrzędnych.
Z definicji koneksji współczynniki pochodne są współrzędnymi
$\nabla_{\partial/\partial t}\omega$. Zależność od bazy
znika, ponieważ parowanie jest naturalne.
\end{proof}

\section{Filtracja Hodge'a i warunek transwersalności Griffithsa}
\begin{definicja}[Wariacja czystych struktur Hodge'a]
Niech $m\in\Z$ i $S$ będzie zespoloną rozmaitością.
\emph{Całkowitą wariacją czystych struktur Hodge'a wagi $m$}
nazywamy lokalny system wolnych grup abelowych
$\mathbb V_{\Z}$ skończonego rzędu oraz malejącą
filtrację $F^\bullet\mathcal H$ przez holomorficzne
podwiązki $\mathcal H=\mathbb V_{\Z}\otimes_{\Z}\mathcal O_S$
o stałych rzędach, spełniającą:
\begin{enumerate}[label=\textup{(\roman*)},leftmargin=*]
\item dla każdego $s\in S$ filtracja na
$\mathbb V_{\Z,s}\otimes\C$ wyznacza czystą strukturę Hodge'a wagi $m$;
\item dla płaskiej koneksji $\nabla$ zachodzi
$\nabla(F^p\mathcal H)\subseteq
F^{p-1}\mathcal H\otimes\Omega_S^1$ dla każdego $p\in\Z$.
\end{enumerate}
Ostatni warunek nazywamy \emph{transwersalnością Griffithsa}.
Polaryzowana wariacja zawiera ponadto płaskie
całkowite parowanie $Q$ na lokalnym systemie,
spełniające $Q(v,w)=(-1)^mQ(w,v)$ oraz
warunki polaryzacji Hodge'a na każdym włóknie
(w szczególności dodatniość odpowiednich form
hermitowskich na składnikach Hodge'a).
\end{definicja}
\begin{uwaga}
Warunek transwersalności nie stwierdza, że
$\nabla F^p\subset F^p$. To ostatnie oznaczałoby
stałość podwiązki Hodge'a względem transportu i jest
ogólnie fałszywe. Ponadto samo pojęcie lokalnego systemu
nie wystarcza do zdefiniowania wariacji Hodge'a: potrzeba
holomorficznej filtracji i jej zgodności z koneksją.
\end{uwaga}
\begin{twierdzenie}[Wariacja Hodge'a z rodziny projektowej]
\label{thm:xv-vhs-family}
Niech $S$ będzie gładką zespoloną odmianą algebraiczną,
a $f:\mathcal X\to S$ gładkim, surjektywnym
morfizmem projektowym o gładkiej dziedzinie $\mathcal X$.
Dla każdego $k\ge0$ lokalny system
$R^kf_*\Z_{\mathcal X}/\mathrm{tors}$ wraz z filtracją
pochodzącą od względnego kompleksu de Rhama
$\Omega^\bullet_{\mathcal X/S}$ tworzy polaryzowaną
wariację czystych struktur Hodge'a wagi $k$.
\end{twierdzenie}
\begin{proof}[Wyodrębnienie części elementarnej i wyniku używanego]
Z twierdzenia Ehresmanna otrzymujemy lokalny system,
a ze Stwierdzenia~\ref{prop:xv-gm-local} płaską koneksję.
Względna filtracja form przez ich stopień holomorficzny
wyznacza podprzestrzenie $F^pH^k(X_s;\C)$.
Stałość ich wymiarów i fakt, że tworzą podwiązki holomorficzne,
wynikają z teorii kohomologii koherentnej i zmiany podstawy
dla gładkich rodzin projektowych. W każdym włóknie
rozkład Hodge'a daje wymagane czyste struktury.

Pozostaje warunek (ii). Lokalny dowód jego postaci
na poziomie kompleksów można uzyskać tak:
wybieramy holomorficzne współrzędne $(t,z_1,\dots,z_d)$
w otoczeniu punktu całkowitej przestrzeni, w których
$f(t,z)=t$. W różniczce absolutnej reprezentanta
względnej formy typu $p$ składniki zawierające $\dd t$
powstają przez różniczkowanie współczynników i przez
porównanie lokalnych podniesień form względnych.
Po wydzieleniu czynnika $\dd t$ mają stopień
względny co najmniej $p-1$. Ten lokalny rachunek dowodzi
właściwego ograniczenia stopnia, lecz jeszcze nie pokazuje,
że otrzymana koneksja jest niezależna od reprezentanta klasy.
Do tego potrzebna jest zgodność różniczki absolutnej
z filtracją względnego kompleksu de Rhama i przejście
przez hiperkohomologię; jest to nietrywialne twierdzenie
porównawcze. Po jego zastosowaniu otrzymuje się
$\nabla F^p\subseteq F^{p-1}\otimes\Omega_S^1$.
Niezależność od wybranych podniesień i sklejanie
tego rachunku na kompleksie \v{C}echa--de Rhama
jest standardowym twierdzeniem o konstrukcji koneksji
Gaussa--Manina; korzystamy z niego w tym miejscu.
Polaryzacja pochodzi z formy przecięcia i klasy
hiperpłaszczyzny przez twierdzenie Hodge'a--Riemanna.
Nie przedstawiamy tego ostatniego jako wyniku
dowiedzionego wyłącznie powyższym rachunkiem.
\end{proof}

\section{Okresy krzywej eliptycznej i półpłaszczyzna okresów}
\begin{stwierdzenie}[Parametr okresów]
Niech $E$ będzie gładką krzywą eliptyczną, $a,b$ zorientowaną
bazą $H_1(E;\Z)$ o $\langle a,b\rangle=1$, a
$0\ne\omega\in H^{1,0}(E)$. Okresy
$A=\int_a\omega$, $B=\int_b\omega$ spełniają $A\ne0$ i
$\tau=B/A\in\mathbb H=\{z\in\C:\operatorname{Im}z>0\}$
przy zgodnym wyborze orientacji formy przecięcia.
\end{stwierdzenie}
\begin{proof}
Gdyby $A=0$, wszystkie okresy formy $\omega$ musiałyby
być rzeczywiście wielokrotnościami $B$, co przeczy
jej dodatniemu polaryzowanemu parowaniu.
Ściślej, relacja Riemanna dla parowania na $H^1(E)$ daje
$\mathrm i\int_E\omega\wedge\overline\omega
=\mathrm i(A\overline B-B\overline A)=2\operatorname{Im}(\overline A B)>0$.
Lewa strona jest dodatnia, ponieważ w lokalnej
współrzędnej $\omega=h(z)\,\dd z$ i
$\mathrm i\,\omega\wedge\overline\omega$
jest dodatnią formą objętościową (z przyjętą orientacją).
Stąd $\operatorname{Im}(\overline AB)>0$, a więc $A\ne0$
i $\operatorname{Im}(B/A)>0$.
\end{proof}
\begin{uwaga}[Zmiana bazy cykli]
Po zamianie bazy $(a,b)$ na inną dodatnio zorientowaną
bazę przez macierz z $\mathrm{SL}_2(\Z)$ parametr
$\tau$ zmienia się ułamkowo liniowo. Nie jest więc
funkcją jednowartościową na dowolnej bazie parametrów:
po obejściu pętli może ulec działaniu monodromii.
\end{uwaga}

\section{Równania Picarda--Fuchsa: od płaskiej koneksji do ODE}
\begin{stwierdzenie}[Równanie spełniane przez okresy]
\label{prop:xv-pf-general}
Niech $S$ będzie gładką nierozkładalną krzywą algebraiczną,
$f:\mathcal X\to S$ gładką rodziną projektową,
$k\ge0$, a
$\mathcal H=H^k_{\mathrm{dR}}(\mathcal X/S)$
algebraiczną wiązką de Rhama rzędu $r\ge1$
z koneksją Gaussa--Manina.
Niech $\omega$ będzie jej przekrojem wymiernym. Ustalmy niestałą funkcję wymierną $t\in\C(S)$
i ograniczmy się do otwartego podzbioru, na którym $t$
jest regularna oraz $\dd t\neq0$. Wtedy $t$ jest
lokalną współrzędną analityczną, a różniczkowanie
$\partial_t$ zachowuje ciało $\C(S)$. Istnieją
funkcje wymierne $a_0,\dots,a_r\in\C(S)$, nie wszystkie zerowe,
takie że każda lokalna funkcja okresowa
$\Pi(t)=\langle\gamma_t,\omega_t\rangle$,
gdzie $\gamma_t$ jest lokalnie poziomą rodziną
cykli w $H_k(X_t;\Z)$, spełnia
$\sum_{j=0}^ra_j(t)\Pi^{(j)}(t)=0$
w dziedzinie, na której współczynniki są określone.
\end{stwierdzenie}
\begin{proof}
Nad ciałem funkcji wymiernych $K=\C(S)$ przestrzeń
generycznego włókna $\mathcal H\otimes K$ ma wymiar $r$.
Ponieważ $t\in K$ jest niestałe, różniczkowanie względem
$t$ jest derywacją ciała $K$ (na krzywej algebraicznej nad
$\C$ mamy $\Omega^1_{K/\C}=K\,\dd t$).
Algebraiczna koneksja przenosi tę derywację na $\mathcal H\otimes K$.
Wektory $\omega,\nabla_{\partial_t}\omega,\dots,
\nabla_{\partial_t}^{r}\omega$ są liniowo zależne nad $K$.
Bierzemy nietrywialną relację z $a_j\in K$.
Parujemy ją z lokalnie poziomym cyklem $\gamma_t$ i
korzystamy $r$ razy ze
Stwierdzenia~\ref{prop:xv-period-derivative}.
Daje to żądane równanie. Algebraiczność koneksji
i przekroju jest w założeniach istotna:
bez niej otrzymalibyśmy współczynniki meromorficzne,
a niekoniecznie wymierne na $S$.
\end{proof}

\section{Rodzina Legendre'a: rachunek okresu od początku}
\label{sec:xv-legendre}
Niech $S=\mathbb P^1\setminus\{0,1,\infty\}$ i rozważmy
rodzinę gładkich krzywych eliptycznych będących
uzupełnieniami rzutowymi krzywych afinicznych
\[
E_\lambda:\qquad y^2=x(x-1)(x-\lambda),\qquad \lambda\in S.
\]
Jej wyróżnik w postaci Weierstrassa wynosi
$\Delta(\lambda)=16\lambda^2(1-\lambda)^2$.
Rzeczywiście wyróżnik wielomianu po prawej stronie,
którego pierwiastki to $0,1,\lambda$, jest
$\lambda^2(1-\lambda)^2$, a wyróżnik równania
Weierstrassa $y^2=P_3(x)$ jest szesnastokrotnością tej liczby.
Forma $\omega_\lambda=\dd x/y$ jest holomorficzna i
nigdzie niezerowa na krzywej rzutowej $E_\lambda$.

\begin{stwierdzenie}[Rzeczywisty okres]
\label{prop:xv-legendre-period}
Niech $E_\lambda$ będzie gładką krzywą eliptyczną
daną rzutowym uzupełnieniem równania
$y^2=x(x-1)(x-\lambda)$ i niech
$\omega_\lambda=\dd x/y$. Dla $0<\lambda<1$
niech $a_\lambda\in H_1(E_\lambda;\Z)$ będzie cyklem,
którego rzut na oś $x$ dwukrotnie przebiega odcinek
$[0,\lambda]$ z odpowiednio wybraną orientacją.
Wówczas okres ma postać
\[
\Pi_a(\lambda)=\int_{a_\lambda}\omega_\lambda
=2\int_0^\lambda\frac{\dd x}
{\sqrt{x(1-x)(\lambda-x)}}
=2\pi\sum_{n=0}^\infty
\frac{\binom{2n}{n}^2}{16^n}\lambda^n.
\]
Szereg zbiega bezwzględnie dla $|\lambda|<1$.
\end{stwierdzenie}
\begin{proof}
Na wnętrzu odcinka $[0,\lambda]$ prawa strona równania
krzywej jest dodatnia, a dwa arkusze nakrycia podwójnego
mają przeciwne znaki $y$. Całka po zamkniętym cyklu
jest więc dwa razy całką po jednym arkuszu.
Podstawienie $x=\lambda u$ daje
$\Pi_a(\lambda)=2\int_0^1
[u(1-u)(1-\lambda u)]^{-1/2}\,\dd u$.
Dla $|\lambda|<1$ rozwijamy
$(1-\lambda u)^{-1/2}=
\sum_{n\ge0}\binom{2n}{n}\lambda^nu^n/4^n$.
Na dowolnym zwartym poddysku $|\lambda|\le r<1$
szereg jest jednostajnie zdominowany przez
stałą pomnożoną przez $[u(1-u)]^{-1/2}$,
więc możemy całkować wyraz po wyrazie.
Po podstawieniu $u=\sin^2\theta$
otrzymujemy
\[
\int_0^1 u^{n-1/2}(1-u)^{-1/2}\,\dd u
=2\int_0^{\pi/2}\sin^{2n}\theta\,\dd\theta.
\]
Całki $I_n=\int_0^{\pi/2}\sin^{2n}\theta\,\dd\theta$
spełniają $I_0=\pi/2$ i
$I_n=(2n-1)I_{n-1}/(2n)$ dla $n\ge1$,
co wynika z całkowania przez części.
Stąd
$\int_0^1 u^{n-1/2}(1-u)^{-1/2}\,\dd u
=\pi\binom{2n}{n}/4^n$.
Po pomnożeniu współczynników otrzymujemy tezę.
\end{proof}

\begin{twierdzenie}[Równanie Picarda--Fuchsa rodziny Legendre'a]
\label{thm:xv-legendre-pf}
Niech $E_\lambda$ będzie rzutowym uzupełnieniem
gładkiej krzywej $y^2=x(x-1)(x-\lambda)$,
$\lambda\in\C\setminus\{0,1\}$, a
$\omega_\lambda=\dd x/y$. Każdy lokalny okres
$\Pi(\lambda)=\int_{\gamma_\lambda}\omega_\lambda$,
gdzie $\gamma_\lambda\in H_1(E_\lambda;\Z)$
zmienia się lokalnie poziomo, spełnia równanie
\[
\lambda(1-\lambda)\Pi''(\lambda)
+(1-2\lambda)\Pi'(\lambda)
-\frac14\Pi(\lambda)=0.
\]
\end{twierdzenie}
\begin{proof}
Niech $a_n=\binom{2n}{n}^2/16^n$.
Bezpośrednio z ilorazu kolejnych współczynników wynika
$(n+1)^2a_{n+1}=(n+\frac12)^2a_n$.
Jeśli $u(\lambda)=\sum_{n\ge0}a_n\lambda^n$,
to współczynnik przy $\lambda^n$ w
$\lambda(1-\lambda)u''+(1-2\lambda)u'-u/4$
jest równy
$(n+1)^2a_{n+1}-(n+\frac12)^2a_n=0$.
Wykazaliśmy więc równanie dla niezerowego okresu
$\Pi_a=2\pi u$. Dla ogólnego okresu odwołujemy się
do własności rzędu dwa koneksji Gaussa--Manina i
do związku równania hipergeometrycznego z rodziną
Legendre'a: drugi okres, dla cyklu nad odcinkiem $[\lambda,1]$,
po podstawieniu $x=1-(1-\lambda)s$ ma postać
$2\pi\mathrm i\,u(1-\lambda)$, z wyborem orientacji cyklu
odpowiadającym znakowi $+\mathrm i$.
Dokładniej, całka po drugim cyklu (z wybraną gałęzią
pierwiastka na $[\lambda,1]$) daje, po podstawieniu
$x=1-(1-\lambda)s$, właśnie funkcję
$2\pi\mathrm i\,u(1-\lambda)$, z dokładnością do
ustalonego znaku wynikającego z orientacji cyklu.
Ponieważ operator
$\lambda(1-\lambda)D^2+(1-2\lambda)D-1/4$
jest niezmienniczy względem podstawienia
$\lambda\mapsto1-\lambda$, obie niezależne
funkcje spełniają to równanie; ich niezależność
wynika z odmiennego zachowania przy $\lambda=0$
(stałość pierwszego okresu i logarytmiczna
rozbieżność drugiego, wykazana poniżej).
Ponieważ $H_1(E_\lambda;\C)$ ma wymiar dwa,
przestrzeń lokalnych funkcji okresowych ma wymiar
co najwyżej dwa. Obie wskazane funkcje okresowe są
liniowo niezależne, zatem tworzą jej bazę.
Każdy pozostały okres jest więc ich lokalną
kombinacją liniową ze stałymi współczynnikami.
\end{proof}
\begin{uwaga}[Kontrola normalizacji]
To jest równanie hipergeometryczne
${}_2F_1(\frac12,\frac12;1;\lambda)$.
Ważny jest znak członu zerowego: stoi przed nim
$-1/4$, a nie $+1/4$. Inna normalizacja okresu
może mnożyć rozwiązanie przez stałą, nie zmienia jednak
samego równania.
\end{uwaga}

\begin{stwierdzenie}[Logarytm i monodromia Legendre'a]
\label{prop:xv-legendre-T2}
Niech $u(\lambda)={}_2F_1(\frac12,\frac12;1;\lambda)$
oraz $v(\lambda)=\mathrm i\,u(1-\lambda)$, początkowo
na $0<\lambda<1$, z kontynuacją analityczną.
W sąsiedztwie $\lambda=0$ zachodzi
$u(\lambda)=1+O(\lambda)$ oraz
$u(1-\lambda)=-\pi^{-1}u(\lambda)\log\lambda+h(\lambda)$
z funkcją holomorficzną $h$ w otoczeniu zera.
Po dodatnim okrążeniu $\lambda=0$ otrzymujemy
$u\mapsto u$ oraz $v\mapsto v+2u$.
W szczególności w bazie rozwiązań $(u,v)$
monodromia działa na kolumnę $(u,v)^t$ przez macierz
$\left(\begin{smallmatrix}1&0\\2&1\end{smallmatrix}\right)$.
Działanie na bazie przestrzeni rozwiązań $(u,v)$
ma macierz transponowaną; przy porównywaniu z
monodromią homologii trzeba ponadto ustalić,
czy stosujemy transport cykli, czy jego odwrotność.
\end{stwierdzenie}
\begin{proof}
Równanie Picarda--Fuchsa ma w $\lambda=0$
regularny punkt osobliwy z podwójnym wykładnikiem
indycjalnym $0$. Metoda Frobeniusa daje drugie
rozwiązanie postaci $c\,u(\lambda)\log\lambda+h(\lambda)$,
gdzie $h$ jest holomorficzne.
Aby wyznaczyć $c$, korzystamy z całki
$u(1-\lambda)=\pi^{-1}\int_0^1
[s(1-s)(1-(1-\lambda)s)]^{-1/2}\,\dd s$.
Dla $s=1-z$ i małego $z$ mianownik ma część dominującą
$\sqrt{z(z+\lambda)}$. Całka
$\int_0^\varepsilon [z(z+\lambda)]^{-1/2}\,\dd z$
ma asymptotykę $-\log\lambda+O(1)$.
Stąd współczynnik wiodący części logarytmicznej wynosi
$c=-1/\pi$, ponieważ $u(0)=1$.
Tożsamość z całym członem $u(\lambda)\log\lambda$
(a nie tylko wyrazem $\log\lambda$) i holomorficzność
reszty $h$ pochodzą z metody Frobeniusa dla regularnego
punktu osobliwego. Samo oszacowanie całki nie dowodziłoby
tej silniejszej postaci.
Przy dodatnim obrocie logarytm zmienia się o $2\pi\mathrm i$,
zatem $u(1-\lambda)\mapsto u(1-\lambda)-2\mathrm i\,u(\lambda)$
i $v\mapsto v+2u$.
\end{proof}
\begin{przyklad}[Efektywny wzór na pochodną ilorazu okresów]
\label{prz:xv-legendre-wronskian}
Dla $0<\lambda<1$ połóżmy, jak wyżej,
$u(\lambda)={}_2F_1(\tfrac12,\tfrac12;1;\lambda)$,
$v(\lambda)=\mathrm i\,u(1-\lambda)$
oraz $\tau(\lambda)=v(\lambda)/u(\lambda)$.
Funkcja $u$ jest dodatnia na tym przedziale,
bo ma dodatnie współczynniki szeregu.
Obie funkcje $u,v$ spełniają
Twierdzenie~\ref{thm:xv-legendre-pf}.
Oznaczmy ich wrońskian przez
$W=u v'-u'v$.
Dzieląc równanie Picarda--Fuchsa przez
$\lambda(1-\lambda)$, otrzymujemy dla obu
rozwiązań równanie $y''+P y'+Q y=0$, gdzie
\[
 P(\lambda)=\frac{1-2\lambda}{\lambda(1-\lambda)}
 =\frac{\dd}{\dd\lambda}
       \log\bigl(\lambda(1-\lambda)\bigr).
\]
Różniczkując wrońskian i wstawiając
równania dla $u$ i $v$, obliczamy
\[
 W'=uv''-u''v
 =-P(uv'-u'v)=-PW.
\]
W konsekwencji
$\lambda(1-\lambda)W(\lambda)=C$
jest stałe na $(0,1)$.
Wyznaczmy ją, zamiast pozostawiać
nieznaną stałą normalizacyjną.
Ze Stwierdzenia~\ref{prop:xv-legendre-T2}
mamy $u(\lambda)=1+O(\lambda)$,
$u'(\lambda)=O(1)$ oraz
\[
 v(\lambda)
 =-\frac{\mathrm i}{\pi}u(\lambda)
          \log\lambda+\mathrm i h(\lambda).
\]
Stąd
$v'(\lambda)=-\mathrm i/(\pi\lambda)
+O(|\log\lambda|+1)$, a więc
$\lim_{\lambda\to0^+}\lambda W(\lambda)
=-\mathrm i/\pi$.
Otrzymujemy dokładny wzór
\[
 \tau'(\lambda)
 =\frac{uv'-u'v}{u^2}
 =-\frac{\mathrm i}
 {\pi\lambda(1-\lambda)u(\lambda)^2}
 \qquad (0<\lambda<1).
\]
W szczególności $\tau'(\lambda)\ne0$,
zatem podprzestrzeń Hodge'a $F^1H^1(E_\lambda)$
rzeczywiście zmienia się z parametrem:
płaskość koneksji Gaussa--Manina
\emph{nie oznacza stałości filtracji Hodge'a}.
Na przedziale rzeczywistym pochodna $\tau'$
jest ujemnie urojona, co zgadza się z
$\tau(\lambda)=\mathrm i\,u(1-\lambda)/u(\lambda)$.
\end{przyklad}

\begin{uwaga}[Dlaczego pojawia się $2$, a nie $1$?]
W standardowym modelu jednokrotnego wygładzenia węzła
$zw=t$ globalna monodromia krzywej jest transwekcją.
Rodzina Legendre'a w parametrze $\lambda$ nie jest tym samym
modelem jednoparametrowym: przy $\lambda=0$
$\Delta(\lambda)$ ma zero rzędu dwa, a odpowiedni
parametr wygładzający lokalnie jest jednostką pomnożoną
przez $\lambda^2$. Dlatego monodromia na integralnej
homologii jest kwadratem transwekcji,
zależnie od orientacji pętli sprzężonym z
$\left(\begin{smallmatrix}1&2\\0&1\end{smallmatrix}\right)$
lub z jego odwrotnością.
Można to sprawdzić bez odwoływania się do samego
wyróżnika. W otoczeniu $x=y=\lambda=0$ wybieramy
holomorficzną gałąź $\sqrt{1-x}$ i wprowadzamy
$Y=y/\sqrt{1-x}$ oraz $X=x-\lambda/2$.
Wtedy równanie rodziny przechodzi dokładnie w
$Y^2+X^2=\lambda^2/4$.
Po podstawieniu $U=Y+\mathrm iX$, $V=Y-\mathrm iX$
otrzymujemy $UV=\lambda^2/4$. Jest to zwykłe
wygładzenie węzła po bazowej zmianie parametru
$t=\lambda^2/4$. Zatem okrążenie $\lambda=0$
odpowiada dwukrotnemu okrążeniu $t=0$.
Sama obecność zwyczajnego węzła we włóknie
centralnym nie wyznacza bez znajomości całkowitej
rodziny wykładnika potęgi transwekcji.
\end{uwaga}

\begin{center}
\begin{tikzpicture}[>=Stealth,scale=0.89,font=\small]
\draw[->] (-2.3,0)--(2.3,0);
\node[below] at (2.00,-0.10) {$\Re\lambda$};
\draw[->] (0,-1.65)--(0,1.65) node[above] {$\Im\lambda$};
\draw[manifoldblue,thick] (0,0) circle (1.05);
\fill[manifoldred] (0,0) circle (2.4pt);
\node[below right] at (0,0) {$0$};
\draw[manifoldgreen,very thick,->] (1.05,0)
 arc[start angle=0,end angle=320,radius=1.05];
\node[above right] at (0.85,0.85) {$\gamma$};
\node at (0,-2.02) {pętla w $\Delta^*$};
\draw[->,thick] (2.55,0)--(4.05,0);
\node[above] at (3.3,0.10) {okresy};
\draw[->] (4.5,0)--(8.35,0);
\node[below] at (8.05,-0.10) {$\Re\tau$};
\draw[->] (5.75,-0.35)--(5.75,2.75) node[above] {$\Im\tau$};
\draw[manifoldblue,densely dotted,thick] (4.5,0)--(8.1,0);
\fill[manifoldblue] (5.55,1.3) circle (2.2pt);
\fill[manifoldgreen] (7.55,1.3) circle (2.2pt);
\draw[manifoldgreen,thick,->] (5.75,1.3)--(7.35,1.3);
\node[below] at (5.55,1.3) {$\tau$};
\node[above] at (7.55,1.3) {$\tau+2$};
\node at (6.35,-1.02) {górna półpłaszczyzna};
\end{tikzpicture}
\par\smallskip\small Schemat: okrążenie $\lambda=0$ w rodzinie Legendre'a
zmienia lokalnie wybraną gałąź stosunku okresów o $2$.
Wykres nie przedstawia globalnie jednowartościowego odwzorowania.
\end{center}

\section{Quasi-unipotencja, filtracja monodromiczna i graniczna struktura Hodge'a}
\begin{definicja}[Unipotencja i logarytm monodromii]
Niech $V_\Q$ będzie skończeniewymiarową przestrzenią
wektorową, a $T\in\operatorname{GL}(V_\Q)$.
Operator $T$ nazywamy \emph{unipotentnym}, jeżeli
$(T-I)^r=0$ dla pewnego $r\ge1$, a \emph{quasi-unipotentnym},
jeżeli pewna jego dodatnia potęga jest unipotentna.
Jeżeli $T$ jest unipotentny, definiujemy nilpotentny
operator $N=\log T:=\sum_{j\ge1}(-1)^{j+1}(T-I)^j/j$.
Suma jest skończona.
\end{definicja}
\begin{twierdzenie}[Twierdzenie o lokalnej monodromii: zakres]
Niech $S$ będzie gładką zespoloną krzywą algebraiczną,
$\overline S$ jej gładkim zwartym uzupełnieniem,
$p\in\overline S\setminus S$, a
$f:\mathcal X\to S$ gładkim projektowym morfizmem.
Wybierzmy lokalny parametr $t$ w punkcie $p$
i mały przebity dysk $\Delta^*\subset S^{\mathrm{an}}$
wokół $p$. Dla każdego $k\ge0$ monodromia lokalnego
systemu $R^kf_*\Q_{\mathcal X}$ wokół $p$
jest quasi-unipotentna: jej wartości własne
są pierwiastkami z jedności. Po odpowiednim
skończonym rozgałęzionym podstawieniu $t=s^m$
monodromia staje się unipotentna.
\end{twierdzenie}
\begin{proof}[Status]
Jest to twierdzenie o monodromii regularnej osobliwości
w geometrycznej wariacji Hodge'a. W ogólności nie wynika
z faktu istnienia płaskiej koneksji; konieczny jest warunek
geometryczności oraz twierdzenie o regularności.
Przyjmujemy je jako wynik zewnętrzny. Rachunek Legendre'a
powyżej dowodzi unipotencji bezpośrednio dla tej rodziny.
\end{proof}

\begin{definicja}[Filtracja monodromiczna w przypadku $N^2=0$]
Niech $V_\Q$ będzie dwuwymiarową przestrzenią
wektorową nad $\Q$, a $N\in\operatorname{End}(V_\Q)$
operatorem takim, że $N\ne0$ i $N^2=0$.
Filtracją wag monodromii \emph{z centrum $1$} nazywamy
filtrację $0=W_{-1}\subset W_0\subset W_1\subset W_2=V_\Q$
zadaną przez $W_0=\operatorname{im}N$ oraz
$W_1=\ker N$. W rozważanym przypadku
$W_0=W_1$ jest jednowymiarowa.
\end{definicja}
\begin{lemat}[Jednoznaczność nietrywialnej filtracji w rzędzie dwa]
\label{lem:xv-weights-rank2}
Niech $V_\Q$ będzie dwuwymiarową przestrzenią,
a $N\in\operatorname{End}(V_\Q)$ operatorem
takim, że $N\ne0$ i $N^2=0$.
Wśród \emph{rosnących, nietrywialnych} filtracji
\[
 0=W_{-1}\subsetneq W_0=W_1
 \subsetneq W_2=V_\Q,
\]
przedłużonych przez $W_j=0$ dla $j<0$
i $W_j=V_\Q$ dla $j\ge2$,
istnieje dokładnie jedna spełniająca
$NW_j\subseteq W_{j-2}$ dla wszystkich $j$
oraz warunek, że
$N:\operatorname{Gr}_2^WV\to\operatorname{Gr}_0^WV$
jest izomorfizmem. Jest nią
$W_0=W_1=\operatorname{im}N=\ker N$.
\end{lemat}
\begin{proof}
Nilpotentny niezerowy operator na przestrzeni
wymiaru dwa ma rząd jeden, a z $N^2=0$ wynika
$\operatorname{im}N\subset\ker N$. Obie przestrzenie
mają wymiar jeden, zatem są równe.
Ścisłe zawierania w hipotezie i warunek
$W_0=W_1$ wymuszają $\dim W_0=\dim W_1=1$.
Ponieważ $NW_2\subseteq W_0$, podprzestrzeń
$W_0$ zawiera jednowymiarowy obraz $N$,
a zatem $W_0=\operatorname{im}N$.
Z drugiej strony $NW_1\subseteq W_{-1}=0$,
więc $W_1\subseteq\ker N$; równość wymiarów
daje $W_1=\ker N$.
Dzięki założeniu \emph{nietrywialności} nie
dopuszczamy sztucznej filtracji z zerowymi
skrajnymi ilorazami, w której odwzorowanie
$0\to0$ byłoby formalnie izomorfizmem.
Odwrotnie, dla $W_0=W_1=\operatorname{im}N=\ker N$
wszystkie zawierania $NW_j\subseteq W_{j-2}$ zachodzą.
Indukowane przez $N$ odwzorowanie
$\operatorname{Gr}_2^WV\to\operatorname{Gr}_0^WV$
jest niezerowym odwzorowaniem między
jednowymiarowymi przestrzeniami, a więc izomorfizmem.
\end{proof}

\begin{przyklad}[Filtracja wag monodromii Legendre'a — pełny rachunek]
\label{prz:xv-legendre-weight-explicit}
Niech $V_\Q=\Q e_0\oplus\Q e_2$.
Wybieramy konwencję transportu, w której
monodromia Legendre'a na dwuwymiarowym
racjonalnym lokalnym systemie ma postać
\[
 T(e_0)=e_0,\qquad T(e_2)=e_2+2e_0,
 \qquad
 [T]_{(e_0,e_2)}
 =\begin{pmatrix}1&2\\0&1\end{pmatrix}.
\]
Ta macierz jest transponowana do macierzy
działania na kolumnie wybranych funkcji okresowych
ze Stwierdzenia~\ref{prop:xv-legendre-T2};
przejście od cykli do funkcji okresowych wymaga
ustalenia działania dualnego oraz orientacji
transportu. Wybrana konwencja wystarcza
do poniższego rachunku, a jej odwrócenie
zmieni jedynie znak operatora $N$.

Bez korzystania z ogólnego twierdzenia o filtracji
obliczamy
\[
 (T-I)^2=0,\qquad
 N=\log T=T-I
 =\begin{pmatrix}0&2\\0&0\end{pmatrix},
 \qquad
 N(e_0)=0,\quad N(e_2)=2e_0.
\]
Stąd
$\operatorname{im}N=\ker N=\Q e_0$
i zgodnie z Lematem~\ref{lem:xv-weights-rank2}
\[
 0=W_{-1}\subset
 W_0=W_1=\Q e_0\subset W_2=V_\Q.
\]
Przestrzeń $\operatorname{Gr}_0^WV$ ma bazę $[e_0]$,
przestrzeń $\operatorname{Gr}_2^WV$ ma bazę $[e_2]$,
a $\operatorname{Gr}_1^WV=0$.
Operator indukowany przez $N$ wysyła
$[e_2]\mapsto2[e_0]$.
Jest to izomorfizm nad $\Q$, dlatego
warunek filtracji monodromicznej
sprawdzamy w tym przykładzie \emph{bezpośrednio}.

Po rozgałęzionym podstawieniu parametru
$\lambda=s^q$, $q\ge1$, dodatnie okrążenie $s=0$
wywołuje $q$ okrążeń $\lambda=0$.
Ponieważ $N^2=0$, a $q$ jest dodatnią liczbą całkowitą,
\[
 T_q=T^q=I+qN
 =\begin{pmatrix}1&2q\\0&1\end{pmatrix},
 \qquad
 N_q=\log T_q=qN.
\]
Mamy $\operatorname{im}N_q=\operatorname{im}N$
i $\ker N_q=\ker N$.
Zatem \emph{filtracja wag pozostaje ta sama},
chociaż całkowity współczynnik monodromii
rośnie z $2$ do $2q$.
Ten wniosek jest rachunkiem liniowym,
nie dowodem istnienia granicznej filtracji
Hodge'a $F_\infty^\bullet$: ta ostatnia
wymaga twierdzenia Schmida.

Porównajmy ten przypadek z Przykładem~\ref{prz:monodromy-zm-v14}.
Dla $z\mapsto z^m$ mamy $T^m=I$, więc po zmianie parametru
$t=s^m$ monodromia na zredukowanej homologii
staje się identycznością. Natomiast w rodzinie Legendre'a
$T^q=I+qN\ne I$ dla \emph{każdego} $q\ge1$,
gdyż pracujemy nad ciałem charakterystyki zero i $N\ne0$.
W obu sytuacjach monodromia jest quasi-unipotentna,
lecz tylko w pierwszej można ją zabić skończonym
rozgałęzionym podstawieniem. Nie utożsamiamy zatem
quasi-unipotencji z unipotencją ani unipotencji
z identycznością.
\end{przyklad}

\begin{figure}[!htbp]
\centering
\begin{tikzpicture}[>=Stealth,font=\small]
 \node[draw,rounded corners=4pt,fill=manifoldblue!8,
       minimum width=3.35cm,minimum height=1.05cm,
       align=center] (a) at (0,0)
       {$\operatorname{Gr}_2^W=\Q[e_2]$\\
        waga $2$};
 \node[draw,rounded corners=4pt,fill=manifoldgreen!8,
       minimum width=3.35cm,minimum height=1.05cm,
       align=center] (b) at (6.1,0)
       {$\operatorname{Gr}_0^W=\Q[e_0]$\\
        waga $0$};
 \draw[->,thick,manifoldred] (a)--(b)
       node[midway,above=8pt] {$N[e_2]=2[e_0]$};
 \node[align=center] at (3.05,-1.45)
       {$\operatorname{Gr}_1^W=0$,\qquad
        $\operatorname{im}N=\ker N=\Q e_0$};
\end{tikzpicture}
\caption{Dwustopniowa filtracja wag monodromii
rodziny Legendre'a. Działanie $N$ łączy
iloraz wagi $2$ z ilorazem wagi $0$;
rozgałęzione podstawienie skaluje $N$,
ale nie zmienia tej filtracji.}
\label{fig:xv-legendre-weight-explicit}
\end{figure}

\begin{twierdzenie}[Graniczna mieszana struktura Hodge'a: wersja używana]
\label{thm:xv-lmhs}
Niech $(\mathbb V_\Z,F^\bullet)$ będzie polaryzowaną
wariacją struktur Hodge'a wagi $k$ nad $\Delta^*$,
pochodzącą od algebraicznej rodziny, i przypuśćmy,
że monodromia $T$ jest unipotentna.
Wybierzmy punkt bazowy $t_0\in\Delta^*$ i połóżmy
$V_\Q=\mathbb V_{\Z,t_0}\otimes_{\Z}\Q$.
Na uniwersalnym nakryciu $t=e^{2\pi\mathrm iz}$
można utożsamić wszystkie włókna $\mathbb V_\Q$
z ustaloną przestrzenią $V_\Q$, a filtrację Hodge'a zapisać jako
$F^\bullet(z)$. Przy konwencji
$F^\bullet(z+1)=T F^\bullet(z)$ granica
\[
F^\bullet_\infty=\lim_{\operatorname{Im}z\to\infty}
\exp(-zN)F^\bullet(z),\qquad N=\log T,
\]
istnieje w przestrzeni flag (po użyciu kanonicznego
przedłużenia) i wraz z filtracją monodromiczną
$W_\bullet(N,k)$ o centrum $k$ tworzy graniczną
mieszaną strukturę Hodge'a. Konwencję indeksów
ustalamy przez warunki
$NW_j\subseteq W_{j-2}$ oraz izomorfizmy
$N^r:\operatorname{Gr}_{k+r}^{W}V_\Q
\to\operatorname{Gr}_{k-r}^{W}V_\Q$
dla wszystkich $r\ge0$; w tym zapisie nie
stosujemy niejednoznacznego symbolu przesunięcia filtracji.
\end{twierdzenie}
\begin{proof}[Status i sens wzoru]
Wartość $F^\bullet(z)$ zależy od gałęzi logarytmu,
a po zastąpieniu $z$ przez $z+1$ jest przekształcana
przez $T=\exp N$. Wtedy
$\exp(-(z+1)N)\exp(N)F^\bullet(z)
=\exp(-zN)F^\bullet(z)$: skorygowana filtracja
jest jednowartościowa. Jest to bezpośrednio sprawdzony
fragment twierdzenia. Jej przedłużenie przez $t=0$,
istnienie granicy i warunki mieszanej struktury Hodge'a
wymagają twierdzenia Schmida o orbicie nilpotentnej
(oraz odpowiedniej teorii przedłużeń). Nie wynikają
z samego rachunku z eksponentą i tutaj korzystamy
z tego głębokiego twierdzenia.
\end{proof}

\begin{przyklad}[Dwie różne struktury przy degeneracji krzywej węzłowej]
\label{ex:xv-central-limit}
Niech gładka krzywa eliptyczna degeneruje do
nieredukowalnej krzywej wymiernej z jednym węzłem,
a lokalna monodromia na $V=H^1(E_t;\Q)$ jest
nieidentycznościowa i unipotentna.
Wtedy $\dim_\Q V=2$, $N\ne0$, $N^2=0$.
Według Lematu~\ref{lem:xv-weights-rank2}
$\operatorname{Gr}_0^WV\cong\Q$,
$\operatorname{Gr}_1^WV=0$,
$\operatorname{Gr}_2^WV\cong\Q(-1)$ jako czyste
struktury Hodge'a (ostatnia identyfikacja wynika z
polaryzacji i twierdzenia o granicznej strukturze).
Tymczasem dla \emph{singularnego włókna centralnego}
$H^1(E_0;\Q)\cong\Q$ ma jedynie wagę $0$,
jak wykazaliśmy w Przykładzie~\ref{ex:nodal-mhs-xiv}.
Zatem $H^1(E_0)$ i graniczna mieszana struktura
Hodge'a gładkiego włókna nie są tym samym obiektem:
nie zgadzają się nawet ich wymiary.
\end{przyklad}
\begin{uwaga}[Tate'owskie przesunięcie wagi]
Symbol $\Q(-1)$ oznacza jednowymiarową czystą
strukturę Hodge'a wagi $2$, typu $(1,1)$.
W przykładzie powyżej $N$ obniża wagę o dwa;
po uwzględnieniu skręcenia Tate'a otrzymujemy
izomorfizm $N:\operatorname{Gr}_2^WV\to
\operatorname{Gr}_0^WV(-1)$.
\end{uwaga}

\section{Granice twierdzeń i dalsza droga do geometrii algebraicznej}
W tym rozdziale w pełni udowodniliśmy istnienie płaskiej
koneksji z lokalnego systemu, wzór na różniczkowanie okresów,
ogólne kryterium istnienia równania Picarda--Fuchsa
przy wyborze algebraicznego parametru lokalnego,
rozwinięcie szeregu okresu Legendre'a, jego równanie i
monodromię oraz lemat liniowo-algebraiczny o filtracji
monodromicznej. Przyjęliśmy natomiast klasyczne twierdzenia
o wariacjach Hodge'a pochodzących z rodzin projektowych,
regularności i quasi-unipotencji monodromii oraz istnieniu
granicznej mieszanej struktury Hodge'a.
Ich niezależne dowody wymagają analizy asymptotycznej
koneksji z osobliwościami regularnymi, teorii modułów
Hodge'a i głębszych twierdzeń algebraiczno-geometrycznych.
Takie rozdzielenie zapobiega kołowemu używaniu
w teorii snopów perwersyjnych i twierdzenia
o rozkładzie.

